\section{Circle measurement and volumes by finite observations}
\label{sec:limits}

This section represents the circumference constant by a specified PR
rational observation recipe. A \emph{rational observation name} is one fixed PR recipe
which, from the numerical input data and a precision $n$, returns rational
endpoints $l_n\le u_n$, with nested intervals and
$u_n-l_n\le2^{-n}$. Equality of two such names means agreement at every
precision, as in the reader construction of Section~\ref{sec:foundation}.
This is a schema for named recipes, not quantification over arbitrary
functions and not a universal evaluator for PR programs. The formulas
below containing $\pi_{\mathrm g}$ abbreviate equalities of these names;
they do not introduce a new algebraic scalar.

The geometric content will be certified by finite enclosures. A disk is
placed between explicit polygons; a solid is placed between finite unions
of polygonal prisms. Their areas and volumes are the exact algebraic
quantities already constructed in the preceding sections. Enclosure
widths tend to zero with stated PR bounds. Thus the same recipes have the
usual geometric interpretation as circle length, disk area, and the
volumes of the displayed solids. No general assertion that arbitrary
bounded sets have an area or volume is used.

\subsection{A polygonal construction of the circumference constant}

Start with the four unit-circle points $(1,0),(0,1),(-1,0),(0,-1)$.
At each stage insert between consecutive vertices $v,w$ the normalized
midpoint $(v+w)/\len{v+w}$. Put
\[
 N_k=4\cdot2^k,\qquad c_0=0,\quad s_0=1,\qquad
 c_{k+1}=\sqrt{\frac{1+c_k}{2}},\quad
 s_{k+1}=\frac{s_k}{2c_{k+1}}.
\]
The two numbers $c_k,s_k$ will be the dot product and positive determinant
of consecutive unit radius vectors. Set
\begin{gather}
 a_k=c_{k+1},\qquad \ell_k=\sqrt{2-2c_k},\qquad
 L_k=N_k\ell_k,\qquad U_k=\frac{L_k}{a_k},\notag\\
 A_k^-=\frac{L_ka_k}{2},\qquad A_k^+=\frac{U_k}{2}.    \tag{L1}
\end{gather}
All these are finite algebraic computations at each numerical stage $k$.

\begin{theorem}[Nested polygonal bounds with explicit error]
\label{lim:polygons}
The construction gives a convex inscribed $N_k$-gon of perimeter $L_k$
and area $A_k^-$. Scaling it about zero by $1/a_k$ gives a polygon
containing the unit disk, of perimeter $U_k$ and area $A_k^+$.
The numerical bounds satisfy
\begin{gather*}
 0<L_k\le L_{k+1}\le U_{k+1}\le U_k\le8,\qquad
 0<A_k^-\le A_{k+1}^-\le A_{k+1}^+\le A_k^+,\\
 U_k-L_k\le4\cdot2^{-k},\qquad
 A_k^+-A_k^-\le2\cdot2^{-k},\qquad
 A_{k+1}^-=L_k/2.                                      \tag{L2}
\end{gather*}
Every finite circle-inscribed polygon whose vertices follow cyclic order
has perimeter at most $U_k$, for every $k$.
\end{theorem}
\begin{proof}
First prove by induction that
$0\le c_k<1$, $s_k>0$, and $c_k^2+s_k^2=1$.
The base case is explicit. The next denominator is positive because
$c_{k+1}^2=(1+c_k)/2\ge1/2$; and
\[
 s_{k+1}^2=\frac{1-c_k^2}{2(1+c_k)}
           =\frac{1-c_k}{2}=1-c_{k+1}^2.
\]
The strict inequalities are preserved. For consecutive unit vectors
$v,w$ with dot product $c_k$ and determinant $s_k$,
$\len{v+w}=2a_k$. Their inserted midpoint $m=(v+w)/(2a_k)$ satisfies
\[
 \dotp vm=\dotp mw=a_k=c_{k+1},\qquad
 \detp vm=\detp mw=s_{k+1}>0.
\]
It lies strictly inside the positive cone spanned by $v,w$.
The four initial cones are the quadrants. Inserting $m$ divides one
cone into two: for a vector $\alpha v+\beta w$ with
$\alpha,\beta\ge0$, compare $\alpha$ and $\beta$ and rewrite it as
either $(\alpha-\beta)v+2a_k\beta m$ or
$2a_k\alpha m+(\beta-\alpha)w$.
Thus the successive cones still make one full cyclic partition of the
plane, and no vertex is repeated.

The midpoint direction $m$ is a unit normal to the chord $vw$, with
$\dotp mv=\dotp mw=a_k$. The other vertices are on the side
$\dotp mx\le a_k$. To verify this last fact algebraically, use the
orthonormal basis $m,\J m$: the chord endpoints have coordinates
$(a_k,\pm\ell_k/2)$. A unit vector with first coordinate greater than
$a_k$ has smaller absolute second coordinate than $\ell_k/2$ and hence
lies strictly in their positive cone. No other vertex lies there.
The chords are therefore the boundary edges of a convex polygon.
The disk is convex, since
\[
 \normsq{(1-t)x+ty}
 =(1-t)\normsq x+t\normsq y-t(1-t)\normsq{y-x}\le1
 \quad(0\le t\le1,\ \normsq x,\normsq y\le1).
\]
Consequently the inscribed polygon lies in the disk. The scaled
polygon has edge equations $\dotp mx=1$, one for each unit chord
normal $m$. To check containment directly, place a disk point in one of
the successive cones, writing it as $x=\alpha v+\beta w$ with
$\alpha,\beta\ge0$. Cauchy--Schwarz gives
$a_k(\alpha+\beta)=\dotp mx\le1$.
Thus it belongs to the triangle with vertices $0,v/a_k,w/a_k$,
and hence to the scaled polygon.

The chord length is $\ell_k$, and the distance of its line from zero is
$a_k$. Summing the $N_k$ triangular fan areas gives $A_k^-=L_ka_k/2$.
Scaling areas by $1/a_k^2$ and lengths by $1/a_k$ gives $A_k^+,U_k$.
For the numerical estimates, the half-angle equations imply
\[
 \ell_{k+1}^2=\frac{\ell_k^2}{2(1+a_k)},\qquad
 L_{k+1}=\frac{L_k}{a_{k+1}},\qquad
 U_{k+1}=\frac{2L_k}{1+a_k}.                         \tag{L3}
\]
Indeed $\ell_k^2=4(1-a_k^2)$ and
$\ell_{k+1}^2=2(1-a_k)$; the second identity then uses
$a_{k+1}^2=(1+a_k)/2$. Since $0<a_k<1$, (L3) gives
$L_{k+1}\ge L_k$ and $U_{k+1}\le U_k$.
Also $U_0=8$, $A_{k+1}^-=L_k/2\ge A_k^-$, and
$A_{k+1}^+=U_{k+1}/2\le A_k^+$.

From $\ell_0^2=2$ and the first equation in (L3), induction gives
$\ell_k^2\le2\cdot2^{-k}$. Now
\[
 U_k-L_k=U_k(1-a_k)
   =\frac{U_k\ell_k^2}{4(1+a_k)}\le2\ell_k^2,
 \qquad
 A_k^+-A_k^-=\frac{U_k(1-a_k^2)}2
            =\frac{U_k\ell_k^2}{8}\le\ell_k^2.
\]
These prove (L2).

For the last assertion, recall the elementary finite perimeter comparison
for convex polygons. Cutting a convex polygon by a half-plane replaces
the removed boundary chain by the straight segment joining its endpoints;
the triangle inequality says that perimeter cannot increase.
If a convex polygon $P$ lies in a convex polygon $Q$, clip $Q$ by the
finitely many supporting half-planes of $P$. The result is $P$, so
$\operatorname{perimeter}(P)\le\operatorname{perimeter}(Q)$.
The circle-inscribed polygon is convex by its cyclic order and is
contained in our outer polygon. Applying this finite comparison proves
the bound by $U_k$.

All vertex invariants in this proof are bounded Boolean folds over the
current vertex list. Induction on $k$ and the finite clipping lists is
therefore quantifier-free induction on coded data. A point whose
membership is being checked is retained as a free parameter.
\end{proof}

\begin{lemma}[Turning algebraic enclosures into rational observations]
\label{lim:rationalize}
Suppose one fixed PR family of algebraic endpoints has proved invariants
$\alpha_k\le\alpha_{k+1}\le\beta_{k+1}\le\beta_k$ and
$\beta_k-\alpha_k\le M2^{-k}$, where $M\ge1$ is an integer computed
from the input data. It has a PR rational observation name with width
at most $2^{-n}$. One explicit choice of algebraic refinement is
\[
 k=n+2+\lambda(M),\qquad
 \lambda(M)=\min\{j\le M:M\le2^j\}.                 \tag{L4}
\]
\end{lemma}
\begin{proof}
The bounded minimum exists because $2^M\ge M$ for $M\ge1$.
Read rational enclosures for $\alpha_k$ and $\beta_k$, each of width
at most $2^{-(n+2)}$. Use the lower endpoint of the first and the upper
endpoint of the second. Their difference is at most
\[
 M2^{-k}+2\cdot2^{-(n+2)}
       \le3\cdot2^{-(n+2)}\le2^{-n}.
\]
To make these rational intervals nested, replace the $n$th interval by
the intersection of the first $n+1$ intervals, taking the maximum lower
endpoint and minimum upper endpoint. It is nonempty: the algebraic
interval at the largest refinement index is contained in all of them.
The last raw width is still an upper bound for the intersection width.
Finite maxima, minima, and the readers from the foundation are PR.
The stated invariants of the intersections are finite rational
inequalities, so their induction is quantifier-free.
This lemma applies to specified families whose invariants have been
proved; it does not take the code of an arbitrary PR function as input.
\end{proof}

\begin{theorem}[Geometric $\pi$, circumference, and disk area]
\label{lim:circle}
The brackets
\[
 p_k=L_k/2\le q_k=U_k/2,
 \qquad 0<q_k-p_k\le2\cdot2^{-k},\qquad q_k\le4       \tag{L5}
\]
define a fixed PR observation name $\pi_{\mathrm g}$.
The unit disk has area $\pi_{\mathrm g}$ and its boundary has
circumference $2\pi_{\mathrm g}$ in the finite-enclosure sense.
For a nonnegative exact algebraic radius $r$,
\[
 \operatorname{Circumference}(r)=2\pi_{\mathrm g}r,
 \qquad \operatorname{DiskArea}(r)=\pi_{\mathrm g}r^2.              \tag{L6}
\]
The circumference brackets $[rL_k,rU_k]$ have width at most
$4r2^{-k}$, and the area brackets $[r^2A_k^-,r^2A_k^+]$ have width at
most $2r^2 2^{-k}$.
\end{theorem}
\begin{proof}
Theorem~\ref{lim:polygons} gives the nested intervals (L5), so apply
Lemma~\ref{lim:rationalize} with $M=2$.
No new exact scalar is being constructed: $\pi_{\mathrm g}$ is the name
of that rational observation recipe.

The polygonal circumference brackets are exactly twice (L5).
Every cyclic polygon inscribed in the circle is bounded above by each
$U_k$, while $L_k$ itself is an inscribed perimeter. Since the gap
$U_k-L_k$ has the displayed modulus, these observations measure the
circle's polygonal length, without a general rectifiability principle.
For area, the actual inner and outer polygons have areas
$A_k^-$ and $A_k^+=q_k$. Their gap has the displayed modulus, and
$A_{k+1}^-=p_k$. Thus the area intervals and (L5) have a common nested
refinement and specify the same observation name.

Multiplying coordinates by $r$ multiplies each finite edge length by
$r$ and each determinant area by $r^2$, proving the enclosures and (L6).
For $r=0$ all measurements are zero. For an arbitrary coded algebraic
$r$, a coarse reader supplies an integer upper bound $B\ge1+r$.
The rationalization lemma can therefore use, for example, $M=4B$ for
circumference and $M=2B^2$ for area. This gives a fully numerical PR
precision schedule, including the cost of rationalizing the endpoints.

When $\pi$ is defined geometrically as half the unit-circle circumference,
$\pi_{\mathrm g}$ is that constant. An independent analytic definition,
such as a trigonometric or integral definition of $\pi$, would require
its own identification theorem. No such identification is assumed in
the computations here.
\end{proof}

For the volume constructions it is convenient to use a matching pair of
disk enclosures: the inscribed polygon at stage $k+1$ has area $p_k$,
and the circumscribed polygon at stage $k$ has area $q_k$.
Their scaled copies enclose every disk of squared radius $b\ge0$ with
areas $p_kb$ and $q_kb$. The required scale $\sqrt b$ is an exact
algebraic operation. The zero-radius case is assigned area zero.

\subsection{Finite slabs for the cylinder, cone, and sphere}

\begin{lemma}[Quadratic slice sums and their modulus]
\label{lim:sums}
For an integer $m\ge1$, define
\[
 S_m^-=\frac{1}{m}\sum_{i=0}^{m-1}\left(\frac{i}{m}\right)^2,
 \qquad
 S_m^+=\frac{1}{m}\sum_{i=1}^{m}\left(\frac{i}{m}\right)^2.
\]
Then
\begin{gather*}
 S_m^-=\frac{1}{3}-\frac{1}{2m}+\frac{1}{6m^2},\qquad
 S_m^+=\frac{1}{3}+\frac{1}{2m}+\frac{1}{6m^2},\\
 0\le S_m^-\le\tfrac13\le S_m^+\le1,\qquad
 S_m^+-S_m^-=\frac{1}{m},                             \tag{L7}\\
 S_m^-\le S_{2m}^-,\qquad S_{2m}^+\le S_m^+.
\end{gather*}
\end{lemma}
\begin{proof}
The integer identity
\[
 6\sum_{i=0}^{m-1}i^2=m(m-1)(2m-1)
\]
is proved by induction: the initial empty sum at $m=0$ is zero, and the
difference of the two right sides at $m+1$ and $m$ is $6m^2$.
This is a quantifier-free polynomial identity; negative factors in the
base expression may be read in the signed-integer encoding.
Adding $m^2$ gives the corresponding identity for $\sum_{i=1}^m i^2$.
Divide by $6m^3>0$ and expand to obtain (L7).
The summands are between zero and one. The lower correction to $1/3$
is $(-3m+1)/(6m^2)\le0$, while the upper correction is positive.
Finally
\[
 S_{2m}^--S_m^-=\frac{2m-1}{8m^2}\ge0,
 \qquad S_m^+-S_{2m}^+=\frac{2m+1}{8m^2}>0.
\]
These calculations prove all the claimed bounds and refinement relations.
No theorem about integration is used.
\end{proof}

A finite slab construction is interpreted as follows. Partition a
vertical interval into finitely many slabs, and on each slab put a fixed
polygonal prism wholly inside the solid and another containing that
part of the solid. Exact prism volumes are base area times height by
Theorem~\ref{sol:prism}. Summing the volumes gives a lower and an upper
bound, since different slabs have disjoint interiors. This sum is a
finite dissection calculation. A common refinement of two slab
partitions verifies comparisons using only finite polygon area
monotonicity. A pair of PR slab bounds with an explicit vanishing gap
therefore provides a volume observation name, without a measure axiom
for arbitrary subsets of space.

\begin{theorem}[Cylinder volume]
\label{lim:cylinder}
For exact algebraic $r,h\ge0$, the cylinder
\[
 x^2+y^2\le r^2,\qquad 0\le z\le h
\]
has the PR volume observation
\[
 V_{\mathrm{cyl}}=\pi_{\mathrm g}r^2h.
\]
At stage $k$ it is enclosed by polygonal prisms of volumes
\[
 C_k^-=p_kr^2h,\qquad C_k^+=q_kr^2h,
 \qquad C_k^+-C_k^-\le2r^2h\,2^{-k}.                 \tag{L8}
\]
\end{theorem}
\begin{proof}
Use the two disk polygons of areas $p_kr^2$ and $q_kr^2$ as bases and
extrude them through height $h$. Disk containment gives containment of
the cylinder between the prisms, and Theorem~\ref{sol:prism} gives
(L8). The numerical bounds are nested by (L5), with the stated gap.
They are exactly $r^2h$ times the brackets defining $\pi_{\mathrm g}$,
so their observation name is the asserted product.
An integer majorant for $2r^2h$ is computed from coarse upper bounds on
$r,h$; Lemma~\ref{lim:rationalize} then supplies rational observations
of any requested width. Zero radius or height gives zero volume directly.
\end{proof}

\begin{theorem}[Cone volume and the one-third ratio]
\label{lim:cone}
For exact algebraic $r\ge0$ and $h>0$, the right circular cone
\[
 0\le z\le h,\qquad x^2+y^2\le r^2(z/h)^2
\]
has the PR volume observation
\[
 V_{\mathrm{cone}}=\frac{\pi_{\mathrm g}r^2h}{3}
                  =\frac{V_{\mathrm{cyl}}}{3}.
\]
With $m=2^k$, finite prism unions give the nested bounds
\[
 K_k^-=p_kr^2hS_m^-,\qquad K_k^+=q_kr^2hS_m^+,
 \qquad K_k^+-K_k^-\le6r^2h\,2^{-k}.                 \tag{L9}
\]
\end{theorem}
\begin{proof}
Divide $[0,h]$ into $m$ slabs of height $h/m$.
On the $i$th slab the cross-section radius lies between
$ri/m$ and $r(i+1)/m$. The disk monotonicity is just the inequality
$i/m\le z/h\le(i+1)/m$ between nonnegative scalars, followed by squaring.
Use the inner polygon of area $p_k(ri/m)^2$ throughout that slab and
the outer polygon of area $q_k(r(i+1)/m)^2$ throughout it.
Their prism unions are finite lower and upper enclosures. Summing
base area times height gives exactly (L9).

Both lower factors $p_k,S^-_{2^k}$ increase, and both upper factors
$q_k,S^+_{2^k}$ decrease; they are nonnegative, so the products are
nested. Their gap is
\begin{align*}
 r^2h(q_kS_m^+-p_kS_m^-)
 &=r^2h\bigl((q_k-p_k)S_m^-+q_k(S_m^+-S_m^-)\bigr)\\
 &\le r^2h\bigl(2\cdot2^{-k}+4/m\bigr)
   =6r^2h\,2^{-k}.
\end{align*}
Here $S_m^-\le1$, $q_k\le4$, and (L5), (L7) supplied all constants.
Moreover the interval
\[
 [p_kr^2h/3,\ q_kr^2h/3]
\]
is contained in $[K_k^-,K_k^+]$, because
$S_m^-\le1/3\le S_m^+$. It is the bracket for
$\pi_{\mathrm g}r^2h/3$. The common refinement and shrinking-width
argument therefore proves equality of the two names. This is the
one-third formula, derived from finite sums rather than a presupposed
integral of a quadratic. An integer majorant for $6r^2h$ makes the
rational precision schedule explicit through (L4).
\end{proof}

\begin{theorem}[Sphere volume and Archimedes' two-thirds ratio]
\label{lim:sphere}
For exact algebraic $R\ge0$, the ball
$x^2+y^2+z^2\le R^2$ has the PR volume observation
\[
 V_{\mathrm{ball}}=\frac{4\pi_{\mathrm g}R^3}{3}.
\]
It is two-thirds of the volume of the cylinder of radius $R$ and height
$2R$. With $m=2^k$, finite prism unions give the nested bounds
\[
 B_k^-=2p_kR^3(1-S_m^+),\qquad
 B_k^+=2q_kR^3(1-S_m^-),\qquad
 B_k^+-B_k^-\le12R^3\,2^{-k}.                        \tag{L10}
\]
\end{theorem}
\begin{proof}
The case $R=0$ has zero volume, so assume $R>0$.
The horizontal section at $z=Rt$, $0\le t\le1$, is a disk of squared
radius $R^2(1-t^2)$. This decreases as $t$ increases.
On the slab $i/m\le t\le(i+1)/m$, take the inner polygon with squared
radius $R^2(1-((i+1)/m)^2)$ and the outer polygon with squared radius
$R^2(1-(i/m)^2)$. The radii themselves are nonnegative algebraic
square roots, so every prism vertex is constructed by the exact scalar
operations. Multiplication by slab height $R/m$ and summation give
the upper-hemisphere bounds
$p_kR^3(1-S_m^+)$ and $q_kR^3(1-S_m^-)$.
Reflect the same slabs through $z=0$ and double their volumes to obtain
(L10).

The lower factors increase and the upper factors decrease, as follows
from (L5), (L7). Put $T^-=1-S_m^+$ and $T^+=1-S_m^-$.
Then $0\le T^-\le2/3\le T^+\le1$ and $T^+-T^-=1/m$.
Exactly as in the cone calculation,
\[
 B_k^+-B_k^-
 =2R^3\bigl((q_k-p_k)T^-+q_k(T^+-T^-)\bigr)
 \le12R^3\,2^{-k}.
\]
The interval $[4p_kR^3/3,4q_kR^3/3]$ lies inside these volume bounds,
so it defines their common observation name. This proves the ball
formula. The circumscribing cylinder has volume
$2\pi_{\mathrm g}R^3$ by Theorem~\ref{lim:cylinder}; multiplying it by
$2/3$ gives the ball volume. Equivalently, the ball and a cone of radius
$R$ and height $2R$ have volumes summing to that cylinder's volume.
The explicit constant $12R^3$, bounded by an integer from a coarse
radius observation, supplies the rational precision schedule (L4).
\end{proof}

\paragraph{What the finite formulas prove internally.}
For each requested $n$, the chosen indices $k$ and $m=2^k$ are finite
numbers. Polygon vertices, slab vertices, determinant volumes, and the
final rational endpoints are produced by specified PR routines.
All containment checks concern algebraic inequalities with a candidate
point as a free parameter. The summation and refinement invariants are
quantifier-free identities on the finite lists. No induction on a
statement universally quantified over every future precision is required.
After proving each precision assertion, universal generalization over
$n$ expresses equality of the observation names.

The constants needed for rational observations are collected here;
each may be replaced by any computable integer majorant at least $1$.
\begin{center}
\begin{tabular}{ll}
\toprule
Observable & Width of its algebraic bracket at stage $k$\\
\midrule
$\pi_{\mathrm g}$ & $2\cdot2^{-k}$\\
Circle circumference & $4r\,2^{-k}$\\
Disk area & $2r^2\,2^{-k}$\\
Cylinder volume & $2r^2h\,2^{-k}$\\
Cone volume & $6r^2h\,2^{-k}$\\
Ball volume & $12R^3\,2^{-k}$\\
\bottomrule
\end{tabular}
\end{center}

\subsection{Why arbitrary real-name predicates need a different treatment}

\begin{proposition}[A uniform obstruction to deciding real-name equality]
\label{lim:equality-boundary}
There is a single PR family of nonnegative rational observation names
$x_e$ for which no total computable procedure decides whether $x_e=0$.
Consequently an exact collinearity or line--circle intersection classifier
for all PR observation names cannot be justified by extending the exact
algebraic sign test to those names.
\end{proposition}
\begin{proof}
Fix a coding of programs, and use bounded simulation of program $e$ on
input $e$. At precision $n$, if its first halt occurred at a stage
$t\le n$, output the singleton interval $[2^{-t},2^{-t}]$;
if no halt has occurred, output $[0,2^{-n}]$.
Bounded simulation and bounded search for the first halting stage are
PR in $(e,n)$. These intervals are nested: a newly found halt has
$t=n+1$ and its singleton lies in the preceding interval, while a
continued failure to halt simply halves the upper endpoint. Their
width is at most $2^{-n}$. Thus this is one particular, uniformly PR
observation recipe, not a universal evaluator of PR functions.

The resulting name is zero exactly when program $e$ does not halt on
input $e$. A total computable zero test $D(e)$ would therefore decide
this halting problem. The contradiction can be given directly:
construct a program which on input $j$ halts if $D(j)$ reports zero,
and otherwise runs forever. Apply it to its own program code.
Either possible answer of the total test contradicts that behavior.
This is the ordinary finite-program diagonal argument, not an assumption
that an infinite search eventually answers.

For geometry, the determinant of $(1,0)$ and $(0,x_e)$ is $x_e$, so an
exact collinearity decision for $(0,0),(1,0),(0,x_e)$ would give the
forbidden zero test. Likewise the line $y=1+x_e$ is tangent to
the unit circle when $x_e=0$ and disjoint from it when $x_e>0$.
An exhaustive classifier on such names would again decide zero.
The conclusion is a uniform computability obstruction. It does not
claim that every individual equality is undecidable, or that continuous
geometry is unavailable. The explicitly proved observation schemata
above remain valid because they require neither that test nor a choice
between exact degeneracy and a positive gap for an arbitrary name.
\end{proof}
